\section{Bài 3. Phát hiện missing token [Cơ bản]}

\subsection{Phân tích \& Thiết kế (M0)}
\begin{itemize}
    \item \textbf{Input:}
    \begin{itemize}
        \item File \texttt{data.csv} (1 cột), mỗi dòng là một token
    \end{itemize}

    \item \textbf{Xử lý:}
    \begin{itemize}
        \item \textbf{Case biên:}
        \begin{itemize}
            \item Token rỗng hoặc chỉ chứa khoảng trắng $\rightarrow$ missing
            \item Token \texttt{"na"}, \texttt{"nan"}, \texttt{"n/a"} (không phân biệt hoa/thường) $\rightarrow$ missing
            \item Token \texttt{"0"} $\rightarrow$ \textbf{không phải missing}
            \item File rỗng $\rightarrow$ không có dòng nào bị missing
        \end{itemize}

        \item \textbf{Module 1 (Chuẩn hóa token):}
        \begin{itemize}
            \item Hàm \texttt{trim(string)}: bỏ khoảng trắng đầu/cuối
            \item Hàm \texttt{toLowerStr(string)}: chuyển về chữ thường
        \end{itemize}

        \item \textbf{Module 2 (Kiểm tra missing):}
        \begin{itemize}
            \item Hàm \texttt{isMissingToken(string t)}
            \item Quy trình:
            \begin{itemize}
                \item trim $\rightarrow$ tolower
                \item nếu rỗng hoặc thuộc \{\texttt{na, nan, n/a}\} $\rightarrow$ true
                \item ngược lại $\rightarrow$ false
            \end{itemize}
        \end{itemize}

        \item \textbf{Module 3 (Đọc file \& phát hiện):}
        \begin{itemize}
            \item Hàm \texttt{findMissingLines(path, vector<int>\&)}
            \item Đọc từng dòng bằng \texttt{getline}
            \item Nếu dòng là missing $\rightarrow$ lưu chỉ số dòng (0-based)
        \end{itemize}

        \item \textbf{Module 4 (Ghi output):}
        \begin{itemize}
            \item Hàm \texttt{writeOutput(path, indices)}
            \item In danh sách chỉ số dòng bị missing
        \end{itemize}
    \end{itemize}

    \item \textbf{Output:}
    \begin{center}
    Danh sách chỉ số dòng (0-based) bị missing, phân tách bằng khoảng trắng
    \end{center}
\end{itemize}

\subsection{Mã nguồn C++11 (Tách module)}
\begin{lstlisting}
    #include <iostream>
    #include <fstream>
    #include <vector>
    #include <string>
    #include <algorithm>

    using namespace std;

    // MODULE: Trim
    string trim(const string& s) {
        size_t start = s.find_first_not_of(" \t\r\n");
        size_t end = s.find_last_not_of(" \t\r\n");

        if (start == string::npos) return "";
        return s.substr(start, end - start + 1);
    }

    // MODULE: To lower
    string toLowerStr(string s) {
        for (char& c : s) {
            c = (char)tolower((unsigned char)c);
        }
        return s;
    }

    // MODULE 1: Check missing token
    bool isMissingToken(string t) {
        t = trim(t);
        t = toLowerStr(t);

        if (t.empty()) return true;
        if (t == "na" || t == "nan" || t == "n/a") return true;

        return false;
    }

    // MODULE 2: Read file & detect missing lines
    bool findMissingLines(string path, vector<int>& indices) {
        ifstream file(path);
        if (!file.is_open()) return false;

        string line;
        int idx = 0;

        while (getline(file, line)) {
            if (isMissingToken(line)) {
                indices.push_back(idx);
            }
            idx++;
        }

        file.close();
        return true;
    }

    // MODULE 3: Write output
    bool writeOutput(string path, const vector<int>& indices) {
        ofstream file(path);
        if (!file.is_open()) return false;

        for (size_t i = 1; i < indices.size(); i++) {
            file << indices[i]+1;
            if (i != indices.size() - 1) file << " ";
        }

        file.close();
        return true;
    }

    // MAIN
    int main() {
        vector<int> indices;

        if (!findMissingLines("data.csv", indices)) {
            cerr << "Khong mo duoc file data.csv\n";
            return 1;
        }

        if (!writeOutput("output.txt", indices)) {
            cerr << "Khong ghi duoc file output.txt\n";
            return 1;
        }

        return 0;
    }\end{lstlisting}

\subsection{Kế hoạch kiểm thử (Test Cases)}

\textbf{Test Case 1 (Thông thường)}

\textbf{Input (\texttt{data.csv}):}
\begin{lstlisting}
10
NA
hello
nan
N/A
0

\end{lstlisting}

\textbf{Expected Output:}
\begin{lstlisting}
2 4 5
\end{lstlisting}

\vspace{0.5cm}

\textbf{Test Case 2 (Không có missing)}

\textbf{Input (\texttt{data.csv}):}
\begin{lstlisting}
1
2
3
abc
0
\end{lstlisting}

\textbf{Expected Output:}
\begin{lstlisting}

\end{lstlisting}

\vspace{0.5cm}

\textbf{Test Case 3 (Case biên -- file rỗng)}

\textbf{Input (\texttt{data.csv}):}
\begin{lstlisting}
\end{lstlisting}

\textbf{Expected Output:}
\begin{lstlisting}

\end{lstlisting}